\documentclass[12pt]{article}
\usepackage{float}
\usepackage{graphicx}
\usepackage{fontspec}
\usepackage{url}
\usepackage{apacite}
\setmainfont{Times New Roman}
\usepackage[letterpaper,top=2cm,bottom=2cm,left=2cm,right=2cm,marginparwidth=1.5cm]{geometry}

\title{Weekly assignment 5: Legal aspects 2}
\author{Truong Khang Tobias Nguyen}
\date{}

\begin{document}

\maketitle

\section{What I learned}
I actually did not expect that we would follow up with another lecture about the legal aspects with Jens, so I was pleasantly surprised by that. 
\newline
That being said, I expected to go deeper about the legislative side of cybersecurity and was hopeful that we might discuss how organizations actually implement security measures and policies that are compliant to legislation.
\newline
Instead, I learned why cybersecurity regulations exist in the first place and why we should even have them; something that I honestly never really put much thought about. In that sense, it was funny when Jens made the joke that Germans just take rules at face value and follow them considering that I am German myself and have actually never bothered to think about it.
\newline
I have also learned that the terms information security and cybersecurity have different definitions since I thought that they would mean the same and thus, always used them interchangeably. 
\newline
Besides that, it was very interesting to learn why organizations should care about regulations where the reasoning goes beyond a simple \textbf{"You will pay a fine if you don't comply"} although I would have liked if we could have gone a bit more into depth instead of adding financial and reputational losses into that discussion.
\newline
I have also learned why authorities themselves regulate cybersecurity as security is a national task and incidents in a smaller module can potentially cascade into even greater consequences in an entire system, e.g. an attack on a critical sector could cripple entire nations.  In that sense, I also got to know the different types of regulations that are in place and what categories they operate such as sector-specific, data protection and general, more agnostic regulations and how they relate to each other.
\newline
Additionally, I think it was especially helpful that Jens showed us how to navigate the landscape of regulations and how to understand legislative jargon.
\newline
I also find it positive that we got introductions to the different categories of regulations, what their scopes are and why they exist.
\newline
What was also different from the last lecture was the amount of concepts taught today as the last lecture was more or less an introduction to the GDPR while this lecture introduced multiple applications and thus, brought more width to the lecture.
\newline
To cap it off, it was also helpful that we discussed previous security incidents and had the case study to get more hands-on on the regulations. 

\section{What I missed}
Although the lecture was well designed, there are still a few points that I felt were missing. Even though I could just recycle the points I made in the last assignment, I will not just repeat them considering that from a conceptual standpoint, they would not have fitted well into the lecture.
\begin{itemize}
\item \textbf{Organizational implementation:} Considering that we were taught why regulations exist and why they are important, I think that we should have also discussed on how organizations actually implement these regulations into their operations in order to be compliant. Although we discussed cases were organizations were \textbf{NOT} compliant, I would have liked to see examples on how a good implementation of regulations would look like.
\item \textbf{Critical perspectives:} Even though I already made that point in the last assignment, I still believe that it is important to question everything. In that sense, I would have like to know about the criticism that exist towards regulations and why, despite existing critisms, regulations exist in their form and how they could be adressed, or if they are even valid in the first place. Considering that regulations (as most things) are not perfect, it is very interesting to know at what point authorities say that they are good enough to roll out and what compromises had to be made. Additionally, I would have also liked to see comparisons between the security performance between heavily-regulated sectors and weakly-regulated sectors and what implications they could have.
\item \textbf{Student discussions:} The last lecture was heavy on discussions between us students with the StudentPulse platform and these kinds of discussions were absent. Although we had the case study at the end, I believe it would have been good if we had smaller discussions at the end of every checkpoint like in the last lecture to apply the knowledge into smaller components that we finish of later instead of the \textbf{"everything at the end"}-approach we had today.
\end{itemize}

\section{What I would've done differently}
Like the last lecture, I have enjoyed this one as well so there is not much that I can add.
\newline
That being said, I would have changed the case study in a way that it is broken down into several smaller pieces that we would compose into one final solution like we did with the StudentPulse case.
\newline
This way, we can apply the knowledge from the checkpoints earlier and may even reason deeper about them, which in turn can improve the final result.
\newline
In that sense, I always feel like there is a kind of pressure when we get bigger cases were the level of reasoning is not as high as if we had smaller tasks first. To make it short, I would have broken the case study in smaller exercises first and then have the final exercise be the composition of their respective solutions.
\newline
Also, as said in the previous assignment, I would have also liked to discuss the other points mentioned but I understand that this is difficult to actually do because of time constraints.

\end{document} 